% Couche Sorties — Databricks official palette (minimalist)
\begin{figure}[H]
\centering
\resizebox{\textwidth}{!}{%
\begin{tikzpicture}[font=\fontfamily{Montserrat-LF}\selectfont, >=Latex]

% === OUTER CONTAINER ===
\draw[rounded corners=8pt, fill=dboatlight, draw=dbnavy, line width=2pt]
  (-0.3,0.5) rectangle (17.3,-9.4);

\node[fill=dbnavy, text=white, font=\bfseries\large,
      rounded corners=4pt, inner sep=9pt, minimum width=17cm]
  at (8.5,0.15) {COUCHE SORTIES \quad
    \normalfont\small\color{white!65}
    3 formats de déploiement — du POC à la production industrielle};

% ============================================================
% SOCLE COMMUN
% ============================================================
\draw[rounded corners=5pt, fill=white, draw=dboatmed, line width=0.8pt, dashed]
  (0.0,-0.6) rectangle (17.1,-2.2);

\node[font=\bfseries\scriptsize, text=dbnavy, anchor=west]
  at (0.3,-1.05) {SOCLE COMMUN};
\node[font=\scriptsize, text=dbnavy!55, anchor=west]
  at (3.5,-1.05) {Partagé par les 3 formats — seule l'encapsulation diffère};

\foreach \bx/\fname in {
  0.3/units/mig\_\{code\}.py,
  5.2/config/,
  7.8/reports/,
  10.4/framework/
}{
  \node[font=\scriptsize\ttfamily, text=dbnavy!70, anchor=west]
    at (\bx,-1.67) {\fname};
}

% ============================================================
% FORMAT 1 — NOTEBOOKS
% ============================================================
\draw[rounded corners=6pt, fill=white, draw=dboatmed, line width=1.2pt]
  (0.0,-2.4) rectangle (5.5,-9.1);

\draw[rounded corners=5pt, fill=dbnavy]
  (0.0,-2.4) rectangle (5.5,-3.25);
\node[text=white, font=\bfseries\normalsize, anchor=west]
  at (0.2,-2.825) {Notebooks};

\foreach \i/\fname in {
  1/run\_all.py,
  2/units/mig\_\{code\}.py,
  3/config/config.json
}{
  \draw[dboatmed, line width=0.4pt]
    (0.2,{-3.5 - (\i-1)*0.78}) -- (5.3,{-3.5 - (\i-1)*0.78});
  \node[font=\scriptsize\ttfamily, text=dbnavy, anchor=west]
    at (0.35,{-3.35 - (\i-1)*0.78}) {\fname};
}

\node[font=\fontsize{7}{8}\selectfont, text=dbnavy!55, italic,
      anchor=west, text width=5.0cm]
  at (0.3,-6.1) {Scripts Python exécutables directement dans un workspace Databricks.};

\draw[dboatmed, line width=0.5pt] (0.2,-7.6) -- (5.3,-7.6);
\node[fill=dboatmed, draw=dboatmed, rounded corners=4pt,
      font=\bfseries\scriptsize, text=dbnavy,
      inner sep=5pt, minimum width=4.8cm, anchor=south]
  at (2.75,-8.95) {Phase~: Démonstration (POC)};

% ============================================================
% FORMAT 2 — PYTHON WHEEL
% ============================================================
\draw[rounded corners=6pt, fill=white, draw=dboatmed, line width=1.2pt]
  (5.8,-2.4) rectangle (11.3,-9.1);

\draw[rounded corners=5pt, fill=dbnavy]
  (5.8,-2.4) rectangle (11.3,-3.25);
\node[text=white, font=\bfseries\normalsize, anchor=west]
  at (6.0,-2.825) {Python Wheel \texttt{.whl}};

\foreach \i/\fname in {
  1/\{pkg\}/\_\_init\_\_.py,
  2/\{pkg\}/runner.py,
  3/pyproject.toml,
  4/README.md
}{
  \draw[dboatmed, line width=0.4pt]
    (6.0,{-3.5 - (\i-1)*0.78}) -- (11.1,{-3.5 - (\i-1)*0.78});
  \node[font=\scriptsize\ttfamily, text=dbnavy, anchor=west]
    at (6.15,{-3.35 - (\i-1)*0.78}) {\fname};
}

\node[font=\fontsize{7}{8}\selectfont, text=dbnavy!55, italic,
      anchor=west, text width=5.0cm]
  at (6.0,-6.4) {Package Python installable via \texttt{pip}. Semver, reproductible, prêt CI/CD.};

\draw[dboatmed, line width=0.5pt] (6.0,-7.6) -- (11.1,-7.6);
\node[fill=dboatmed, draw=dboatmed, rounded corners=4pt,
      font=\bfseries\scriptsize, text=dbnavy,
      inner sep=5pt, minimum width=4.8cm, anchor=south]
  at (8.55,-8.95) {Phase~: Production};

% ============================================================
% FORMAT 3 — ASSET BUNDLE  (lava accent — flagship DevOps format)
% ============================================================
\draw[rounded corners=6pt, fill=dblava!6, draw=dblava, line width=2pt]
  (11.6,-2.4) rectangle (17.1,-9.1);

\draw[rounded corners=5pt, fill=dblava]
  (11.6,-2.4) rectangle (17.1,-3.25);
\node[text=white, font=\bfseries\normalsize, anchor=west]
  at (11.8,-2.825) {Asset Bundle (DAB)};

\foreach \i/\fname in {
  1/databricks.yml,
  2/bitbucket-pipelines.yml,
  3/scripts/fetch\_reports.ps1,
  4/+ tout le contenu Wheel
}{
  \draw[dblava!20, line width=0.4pt]
    (11.8,{-3.5 - (\i-1)*0.78}) -- (16.9,{-3.5 - (\i-1)*0.78});
  \node[font=\scriptsize\ttfamily, text=dbnavy, anchor=west]
    at (11.95,{-3.35 - (\i-1)*0.78}) {\fname};
}

\node[font=\fontsize{7}{8}\selectfont, text=dbnavy!55, italic,
      anchor=west, text width=5.0cm]
  at (11.8,-6.4) {Bundle complet avec orchestration CI/CD Bitbucket et déploiement Databricks CLI.};

\draw[dblava!25, line width=0.5pt] (11.8,-7.6) -- (16.9,-7.6);
\node[fill=dblava, draw=dblava, rounded corners=4pt,
      font=\bfseries\scriptsize, text=white,
      inner sep=5pt, minimum width=4.8cm, anchor=south]
  at (14.35,-8.95) {Phase~: DevOps (CI/CD)};

\end{tikzpicture}%
}
\caption{Couche Sorties~: trois formats de déploiement générés par le compilateur}
\label{fig:archi-sorties}
\end{figure}
